\section{Related Work}
% FACT SOURCE: notes/report_sources.md Sec. 2. Every citation below was re-checked against the full PDF text
% (not just the abstract) in ref_dump/ on 2026-09-27. Restructured 2026-09-27: each of the three paragraphs now
% opens with a \paragraph{} label naming its sub-question, and its very next sentence states the gap in plain
% language, before the supporting evidence. Budget ~0.85 page incl. table.

\paragraph{SQ1: grading damage for repair.} No prior study grades damage severity for the repair step itself; graded
severity has only been used for \emph{reading} already-OCR'd text. Post-OCR research splits the task into error
\emph{detection} and error \emph{correction} \citep{nguyen2021survey}, scored separately by the ICDAR 2019
competition \citep{rigaud2019icdar}; we study correction only and assume the damaged span is known, the assumption
Method~\ref{sec:method} builds on. Results on how well GenAI corrects disagree: a neural sequence model corrects OCR
errors directly \citep{lyu2021neural}; a Llama~2 model fine-tuned on BLN600 cut CER further than a fine-tuned BART
\citep{thomas2024leveraging}; 14 prompted LLMs were weak across several historical benchmarks
\citep{boros2024postcorrection}; a seq2seq model beat generative alternatives on early modern Dutch
\citep{debaene2025evaluating}; and prompted LLMs lowered CER for English text but not for Finnish
\citep{kanerva2025ocr}. These results diverge along three lines -- fine-tuned versus prompted, the language involved,
and how damaged the source corpus already was -- but every study reports CER or WER averaged over its corpus's
natural error distribution \citep{nguyen2021survey}, which cannot show at what point repair stops working. Graded
severity levels do exist, but only for reading OCR text: synthetic degradations at several severities have measured
named-entity recognition \citep{hamdi2020assessing}, and real OCR output has been binned into quality bands for
similar downstream tests \citep{vanstrien2020assessing}. Applying that grading logic to the repair step itself is
what SQ1 does.

\paragraph{SQ2: restoring text at known gaps.} The closest prior task is restoring damaged ancient text at known
positions, and it differs from ours in exactly the ways SQ2 probes. Pythia hides 1--10 characters at known positions
in Greek inscriptions and predicts them \citep{assael2019restoring}; Ithaca extends this and pairs predictions with
historians' own judgement \citep{assael2022restoring}; a masked-language-model approach fills gaps in Akkadian
cuneiform text, with accuracy falling as the masked span grows \citep{lazar2021filling}. Three differences motivate
our design. First, these studies hide characters completely, whereas OCR damage leaves letters garbled yet partly
readable, which our slot input exploits (Sec.~\ref{sec:slot}). Second, their gaps are short, at most ten characters
or a few tokens, while ours reach 75\% of a sentence. Third, none of them compares a masked encoder against a
prompted LLM, or scores facts separately from overall similarity -- the comparison SQ2 makes. On the model side,
BERT fills masked positions in parallel \citep{devlin2019bert}, which makes multi-token blanks harder to fill well
\citep{shen2020blank} and motivates the short beam search over candidate lengths in Method~\ref{sec:models};
left-to-right generative infilling \citep{donahue2020enabling} is the natural counterpart for a prompted LLM; and
combining context with letter-level evidence is a classic OCR correction idea \citep{evershed2014correcting}, of
which our reranking step is a direct, modern instance. This also settles whether a masked encoder counts as GenAI at
all: BERT can be sampled to produce text \citep{wang2019bert}, and filling a blank from context is generation
conditioned on the rest of the sentence \citep{donahue2020enabling}.\footnote{We cite \citet{wang2019bert} only for
this narrow claim; the paper's own Markov-random-field derivation was later retracted by one of its authors, see
\url{https://kyunghyuncho.me/bert-has-a-mouth-and-must-speak-but-it-is-not-an-mrf/}. A reader unconvinced by the
generative framing can instead read BERT as the fine-tuned, in-domain baseline the few-shot LLM has to beat, which
is how our experiment treats it either way.}

\paragraph{SQ3: meaning versus facts.} A similarity score can pass a repair that gets the facts wrong, and SQ3 tests
whether that happens here. BERTScore \citep{zhang2020bertscore} is the standard similarity metric for judging
generated text against a reference. In summarisation it correlates only weakly with human judgements of
faithfulness (Spearman 0.19) and factuality (Spearman 0.12) \citep{maynez2020faithfulness}, and it can give a high
score to a candidate that shares most content words with the reference while stating a different fact
\citep{hanna2021finegrained} -- exactly the failure mode of a repair with the right context and one wrong name. This
evidence comes from summarisation and machine translation, not from text repair, so SQ3 checks it directly by
comparing BERTScore against the Fact Recovery Rate on the same repaired sentences. A second concern is memorisation:
BLN600 has been publicly available since 2024, so an instruction-tuned LLM may have seen parts of it during
pretraining \citep{sainz2023contamination}. We probe for this with guided completion: give the model the first part
of a test instance and compare its continuation with the true remainder \citep{golchin2024time}.

No single element of our design is new by itself; the combination is. Table~\ref{tab:position} places our study
against the three closest lines of work on five points: whether the damage location is known, whether the damaged
letters stay visible, whether severity is graded, whether facts are scored separately from overall similarity, and
whether an encoder and an LLM are compared under the same input. The gaps this leaves are what Method
\ref{sec:method} builds on: graded nested damage (Sec.~\ref{sec:damage}), a shared slot input for both repair
methods (Sec.~\ref{sec:slot}), two repair families evaluated side by side (Sec.~\ref{sec:models}), and a fact-level
score reported next to BERTScore (Sec.~\ref{sec:metrics}).

\begin{table}[t]
\centering
\small
\setlength{\tabcolsep}{2.5pt}
\begin{tabular}{@{}lccccc@{}}
\toprule
 & Known & Letters & Graded & Fact & Enc.\ vs \\
Work & loc. & visible & severity & score & LLM \\
\midrule
Post-OCR LLMs$^a$ & $\times$ & \checkmark & $\times$ & $\times$ & $\times$ \\
Pythia$^b$ & \checkmark & $\times$ & $\times$ & $\times$ & $\times$ \\
Akkadian MLM$^c$ & \checkmark & $\times$ & \checkmark & $\times$ & $\times$ \\
\textbf{Ours} & \checkmark & \checkmark & \checkmark & \checkmark & \checkmark \\
\bottomrule
\end{tabular}
\caption{Position of this study. $^a$\citet{thomas2024leveraging,boros2024postcorrection,kanerva2025ocr}
(real OCR, CER/WER over the corpus). $^b$\citet{assael2019restoring} (1--10 hidden characters).
$^c$\citet{lazar2021filling} (masked tokens, results by span length). Graded: results reported by amount
of damage.}
\label{tab:position}
\end{table}

% ---- NOTES FOR DEFENSE / APPENDIX (kept out of the page count discussion) ----
% - Belinkov and Bisk (2018), synthetic vs natural noise transfer: moved to Limitations, not cited here,
%   since it answers a validity threat (Sec. 6) rather than one of SQ1-SQ3.
% - Hamdi et al. appear twice on purpose: the 2020 paper documents the synthetic graded degradation used
%   here; the 2023 NLE journal version (cited in the Introduction) gives the in-depth NER/NEL analysis.
%   Keep both; they are not duplicates.
% - Restructured 2026-09-27: added \paragraph{} labels (SQ1/SQ2/SQ3) per user request, and moved each
%   paragraph's gap statement to its second sentence so the point is explicit up front, not buried at the
%   end of an unbroken block.
